\section{Estimación del Costo de Software}

La valoración del software construido se apoya en el \textbf{análisis de puntos
de función} \cite{albrecht1979measuring,ifpug2010counting}, método estandarizado que mide el tamaño de un sistema a partir de las
funciones que ofrece al usuario y no de su implementación. Esta propiedad es la
que lo hace adecuado para el presente caso: permite valorar el sistema a partir
de los requerimientos y de la arquitectura documentados en los capítulos
precedentes, con independencia del lenguaje empleado y del estilo de codificación
del desarrollador.

Se descarta explícitamente el uso de modelos paramétricos de estimación
\cite{boehm1981software}
construidos para proyectos industriales. Esos modelos se calibran sobre equipos de
desarrollo, ciclos de vida prolongados y procesos organizacionales formalizados,
condiciones que no se verifican en un proyecto de grado de desarrollo individual.
Aplicarlos produciría una sobreestimación sin sustento; en su lugar se emplea un
método de medición funcional cuyo resultado es verificable.

\subsection{Tamaño funcional del sistema}

El recuento clasifica cada función en una de cinco categorías ---archivos lógicos
internos, archivos de interfaz externa, entradas externas, salidas externas y
consultas externas--- y le asigna una complejidad según la cantidad de tipos de
registro y de elementos de dato que manipula. La Tabla \ref{tab:puntos_funcion}
presenta el recuento resultante.

\begin{longtable}{
        P{0.10}
        P{0.55}
        P{0.16}
        P{0.09}
    }
    \caption{Recuento de puntos de función sin ajustar}
    \label{tab:puntos_funcion}\\

    \toprule
    \rowcolor{gray!20}
    \textbf{Tipo} &
    \textbf{Función} &
    \textbf{Complejidad} &
    \textbf{Puntos} \\
    \midrule
    \endfirsthead

    \multicolumn{4}{c}{\textit{Continuación de la Tabla
    \ref{tab:puntos_funcion}}} \\
    \toprule
    \rowcolor{gray!20}
    \textbf{Tipo} &
    \textbf{Función} &
    \textbf{Complejidad} &
    \textbf{Puntos} \\
    \midrule
    \endhead

    Archivo lógico & Vehículos observados: atributos identificados y matrícula & Baja & 7 \\
    \midrule
    Archivo lógico & Sesiones de estacionamiento: ingreso, salida y permanencia & Media & 10 \\
    \midrule
    Archivo lógico & Detecciones: bitácora de inferencia & Media & 10 \\
    \midrule
    Entrada externa & Registro de ingreso en el punto de acceso & Media & 4 \\
    \midrule
    Entrada externa & Registro de salida en el punto de acceso & Baja & 3 \\
    \midrule
    Entrada externa & Procesamiento del video archivado por el extractor & Media & 4 \\
    \midrule
    Salida externa & Agregados diarios de ocupación y permanencia & Media & 5 \\
    \midrule
    Salida externa & Notificación de variaciones de ocupación & Baja & 4 \\
    \midrule
    Salida externa & Generación del conjunto de datos anotado & Media & 5 \\
    \midrule
    Consulta externa & Sesiones actualmente activas & Baja & 3 \\
    \midrule
    Consulta externa & Historial de sesiones con paginación & Media & 4 \\
    \midrule
    Consulta externa & Nivel de ocupación del recinto & Baja & 3 \\
    \midrule
    \multicolumn{3}{r}{\textbf{Total sin ajustar}} & \textbf{62} \\*
    \multicolumn{4}{r}{\footnotesize\textit{Fuente: Elaboración propia}} \\*
    \bottomrule
\end{longtable}

La complejidad se asignó según el número de tipos de registro y de elementos de
dato que cada función manipula. Se clasificaron como de complejidad baja las
funciones que operan sobre una única estructura de datos con un número reducido
de campos, y como de complejidad media aquellas que involucran una relación entre
dos entidades o que devuelven información derivada de un cálculo agregado. No se
clasificó ninguna función como de complejidad alta, dado que la estructura de
datos del sistema es acotada y no contiene relaciones de tercer nivel ni
agrupaciones jerárquicas complejas.

Al recuento sin ajustar se aplica el factor de ajuste que pondera catorce
características generales del sistema. Su valor se adopta como la unidad por las
siguientes razones: las características que incrementarían el factor ---exigencia
de desempeño, procesamiento complejo derivado del tratamiento de imágenes, y
facilidad de operación e instalación--- se compensan con aquellas que lo
reducirían ---ausencia de procesamiento distribuido, de operación en múltiples
sitios y de necesidad de integración con sistemas externos---. El tamaño funcional
ajustado resulta, en consecuencia, de \textbf{62 puntos de función}.

\subsection{Contraste con el tamaño medido del código}

Una verificación independiente del recuento anterior se obtiene midiendo el
código fuente efectivamente escrito. La Tabla \ref{tab:lineas_codigo} presenta el
resultado del conteo por área.

\begin{longtable}{
        P{0.64}
        P{0.14}
        P{0.14}
    }
    \caption{Tamaño del código fuente por área}
    \label{tab:lineas_codigo}\\

    \toprule
    \rowcolor{gray!20}
    \textbf{Área} &
    \textbf{Archivos} &
    \textbf{Líneas} \\
    \midrule
    \endfirsthead

    \multicolumn{3}{c}{\textit{Continuación de la Tabla
    \ref{tab:lineas_codigo}}} \\
    \toprule
    \rowcolor{gray!20}
    \textbf{Área} &
    \textbf{Archivos} &
    \textbf{Líneas} \\
    \midrule
    \endhead

    Dominio (reglas de negocio y puertos) & 10 & 615 \\
    \midrule
    Aplicación (casos de uso) & 2 & 182 \\
    \midrule
    Adaptadores (base de datos, cola, bus y modelos) & 18 & 1,234 \\
    \midrule
    Entrada HTTP (rutas, esquemas y configuración) & 18 & 1,408 \\
    \midrule
    Trabajador de inferencia & 5 & 349 \\
    \midrule
    Migraciones del esquema & 4 & 236 \\
    \midrule
    Captura de flujos de video & 2 & 361 \\
    \midrule
    Cliente web & 29 & 2,561 \\
    \midrule
    \textbf{Subtotal del producto} & \textbf{88} & \textbf{6,946} \\
    \midrule
    Pruebas del servicio y de las herramientas & 40 & 3,537 \\
    \midrule
    Herramientas de datos y entrenamiento & 31 & 6,225 \\
    \midrule
    \textbf{Total entregado} & \textbf{159} & \textbf{16,708} \\*
    \multicolumn{3}{r}{\footnotesize\textit{Fuente: Elaboración propia}} \\*
    \bottomrule
\end{longtable}

El código se mide sobre los archivos fuente efectivamente entregados, excluyendo
las dependencias de terceros, los artefactos generados y los archivos de
inicialización de paquete, que no contienen lógica. El conteo distingue dos
naturalezas: el \textbf{producto}, que implementa las funciones del sistema, y el
material de \textbf{verificación y de construcción de datos}, que sostiene el
trabajo pero no forma parte de la interfaz ofrecida al usuario.

Ambas mediciones son independientes entre sí: el tamaño funcional se derivó de
las funciones documentadas y el tamaño del código, de los archivos fuente. La
relación obtenida sobre el producto ---$6{,}946$ líneas para $62$ puntos de
función, es decir aproximadamente $112$ líneas por punto de función--- constituye
una verificación de consistencia interna del análisis. Las pruebas y las
herramientas de datos se informan por transparencia, pero se excluyen de esa
relación: contarlas elevaría el cociente a cerca de $270$ líneas por punto de
función y describiría el esfuerzo de construcción del conjunto de datos, no el
tamaño funcional del sistema entregado.

\subsection{Esfuerzo y costo de desarrollo}

El esfuerzo de desarrollo se obtiene a partir del tamaño funcional y de una
productividad de referencia, y el costo, a partir del esfuerzo y de una tarifa:

\begin{equation}
    E = TFA \cdot P \qquad\qquad C_{\text{software}} = E \cdot T
\end{equation}

donde $TFA$ es el tamaño funcional ajustado ---$62$ puntos de función---, $P$ la
productividad expresada en horas por punto de función, $E$ el esfuerzo en horas,
$C_{\text{software}}$ el costo de desarrollo y $T$ la tarifa horaria.

Los dos parámetros de valoración ---productividad y tarifa--- dependen de
condiciones locales de mercado y de la experiencia del desarrollador, y no son
derivables del análisis técnico. Se declaran de forma explícita en la Tabla
\ref{tab:parametros_costo} para que la estimación pueda reproducirse y
recalcularse.

\begin{longtable}{
        P{0.30}
        P{0.12}
        P{0.12}
        P{0.36}
    }
    \caption{Parámetros de valoración del costo de software}
    \label{tab:parametros_costo}\\

    \toprule
    \rowcolor{gray!20}
    \textbf{Parámetro} &
    \textbf{Símbolo} &
    \textbf{Valor} &
    \textbf{Procedencia} \\
    \midrule
    \endfirsthead

    \multicolumn{4}{c}{\textit{Continuación de la Tabla
    \ref{tab:parametros_costo}}} \\
    \toprule
    \rowcolor{gray!20}
    \textbf{Parámetro} &
    \textbf{Símbolo} &
    \textbf{Valor} &
    \textbf{Procedencia} \\
    \midrule
    \endhead

    Tamaño funcional ajustado & $TFA$ & 62 puntos & Calculado en el apartado anterior. \\
    \midrule
    Productividad & $P$ & $12$ horas por punto & Parámetro declarado. \\
    \midrule
    Esfuerzo & $E$ & $744$ horas & $TFA \cdot P$. \\
    \midrule
    Tarifa horaria & $T$ & Bs $30$ & Parámetro declarado. \\
    \midrule
    Costo de desarrollo & $C_{\text{software}}$ & Bs 22,320 & $E \cdot T$. \\*
    \multicolumn{4}{r}{\footnotesize\textit{Fuente: Elaboración propia}} \\*
    \bottomrule
\end{longtable}

Los valores adoptados se declaran de la siguiente forma. La productividad se fija
en $12$ horas por punto de función, equivalentes a $744$ horas de dedicación para
el tamaño calculado; el valor se adopta por analogía con proyectos de magnitud
comparable desarrollados en lenguajes de alto nivel, y es coherente con un
desarrollo individual de varios meses. La tarifa se fija en Bs $30$ por hora, que
corresponde a un salario mensual de referencia de Bs 4,800 sobre una jornada de
$160$ horas. Con ambos parámetros, el costo de desarrollo resulta de Bs 22,320.

La estimación es sensible a la tarifa y no al tamaño: con una tarifa de Bs $20$
el costo de software desciende a Bs 14,880, y con Bs $40$ asciende a Bs 29,760,
mientras que el tamaño funcional permanece fijo porque se deriva de las
funciones documentadas y no del esfuerzo invertido. Los dos parámetros son de
mercado y deben reemplazarse por la cotización vigente al momento de la
presentación; la estructura del cálculo no cambia al hacerlo.
